\chapter{The Seven-Line Schema and Its Three Amendments}\label{ch:seven-lines}

This chapter prints the schema and explains why the columns have one-letter names. That is all it does. The views and functions are Chapter 8; the identity defect is Chapter 11. Here there are two tables, and a lineage.

\section{The tables, as the database reports them}

Ask the live instance what it has and it answers in fourteen lines, of which seven carry column definitions. This is \texttt{\textbackslash d o} and \texttt{\textbackslash d r}, transcribed from a session on 2026-07-17, after the big bang:

\begin{lstlisting}
CREATE TABLE o (
  i UUID PRIMARY KEY DEFAULT uuidv7(),
  j JSONB NOT NULL
);

CREATE TABLE r (
  i  UUID NOT NULL DEFAULT uuidv7(),
  o1 UUID,
  o2 UUID,
  d  JSONB NOT NULL,
  c  REAL NOT NULL DEFAULT 1,
  b  UUID NOT NULL DEFAULT gen_random_uuid(),
  t  TIMESTAMPTZ NOT NULL DEFAULT now()
);
\end{lstlisting}

Both \texttt{i} columns were \texttt{BIGINT GENERATED BY DEFAULT AS IDENTITY} until 17 July 2026, and every earlier printing of this chapter showed them that way. The column count did not move and the meaning of the columns did not move; what moved is who is allowed to mint a name. Chapter 11 owns that argument. It is mentioned here only so that the seven lines are the seven lines that are actually running.

\texttt{o} holds objects: observations, sources, boundary references, rules, and rules about rules. Its entire content is a JSONB blob \texttt{j}. \texttt{r} holds relations: \texttt{o1} and \texttt{o2} are the endpoints, \texttt{d} is the payload (whose \texttt{d->>'type'} is the relation's type), \texttt{c} is confidence, \texttt{b} is a branch, \texttt{t} is when the row was written, and \texttt{i} is the relation's own identity, which is what makes it possible to say something about a relation rather than only with one.

\texttt{r} is a TimescaleDB hypertable partitioned on \texttt{t}:

\begin{lstlisting}
SELECT create_hypertable('r', 't');
ALTER TABLE r SET (
  timescaledb.compress,
  timescaledb.compress_segmentby = 'b',
  timescaledb.compress_orderby   = 't DESC'
);
SELECT add_compression_policy('r', INTERVAL '30 days');
\end{lstlisting}

\begin{Correspondence}{The live schema is two tables and seven columns}{Ask a live instance to describe both tables and compare what it prints against the block above.}
There is no third table. There is no table of persons, no table of samples, no table of papers, no table of outcomes. Everything that a normalized design would give a table gets a \texttt{j} in \texttt{o} and some rows in \texttt{r}.
\begin{Check}
docker exec yoneda-cdb psql -U postgres -d yoneda -c "\d o" -c "\d r"
-- both i columns: uuid, not null, default uuidv7()
-- o reports a primary key on i; r reports none, see the last section
-- r reports as a hypertable on t, with no child table until you insert a row
\end{Check}
\end{Correspondence}

\subsection{The checked-in schema was stale, and then the drift was cancelled rather than fixed}

Before going further, one honest note that belongs in the chapter that prints the code, and it has to be told in the past tense now, which is itself the point.

Until 17 July 2026 the file \texttt{schema.sql} was not a faithful description of the running database, and the drift ran in the direction of live-ahead-of-file. \texttt{schema.sql} created four indexes on \texttt{r}: \texttt{r\_i\_idx}, \texttt{r\_o1o2\_idx}, a GIN index on \texttt{d}, and \texttt{r\_b\_idx}. The live table had eight. The four extra were \texttt{r\_o1\_idx}, \texttt{r\_o2\_idx}, \texttt{r\_t\_idx}, and a partial expression index \texttt{r\_outcome\_of\_idx} on \texttt{((d->>'of')::bigint) WHERE d->>'type' = 'outcome'}, which is exactly the lookup \texttt{r\_active} performs once per row and which nobody wrote down. Likewise \texttt{schema.sql} created one extension; the live database reported five: \texttt{plpgsql}, \texttt{vector}, \texttt{pg\_trgm}, \texttt{timescaledb}, \texttt{timescaledb\_toolkit}.

Every one of those numbers is now wrong, and none of them was corrected. The live instance was rebooted from \texttt{stem-cell/stem-cell.sql}, so the file and the database agree by construction: the database is nothing but what that file created, four minutes of it. The extensions are three, \texttt{plpgsql}, \texttt{timescaledb}, and \texttt{timescaledb\_toolkit}, of which the file asks for exactly one and the image supplies the other two. The indexes on \texttt{r} are five: the four the file creates plus \texttt{r\_t\_idx}, which TimescaleDB creates for you when you call \texttt{create\_hypertable}. \texttt{pg\_trgm} and \texttt{vector} are gone, along with the four hand-added indexes, because the instance that had them is a file on a disk. \texttt{schema.sql} is still in the repository and still describes the \texttt{BIGINT} universe; it is no longer the schema of anything that runs.

\begin{Correspondence}{The file and the database agree, because the database was rebooted from the file}{Count the indexes and extensions a live instance reports, then count what \texttt{stem-cell.sql} creates.}
Do not read this as a repair. The drift was real for a year and nobody closed it; what closed it was wiping the database, which is a remedy available exactly once, at birth, and Chapter 12 is blunt about the price. The general lesson survives its own instance being deleted, and it is the one the system's own rules already state: a file is a snapshot of a mutable thing, dated at the moment it was written, and it stays true only about that moment. This file is true about this morning. It will start drifting the first time somebody adds an index at 2am, and nothing in the design will notice. Chapter 8 finds a view in production that was in no file at all.
\begin{Check}
select extname from pg_extension order by 1;
-- plpgsql, timescaledb, timescaledb_toolkit                     (3)
select indexname from pg_indexes where tablename='r' order by 1;
-- r_b_idx, r_d_idx, r_i_idx, r_o1o2_idx, r_t_idx                (5)
grep -c "CREATE INDEX" stem-cell.sql      -- 4; r_t_idx is Timescale's, not the file's
grep -c "CREATE EXTENSION" stem-cell.sql  -- 1; the image supplies the other two
\end{Check}
\end{Correspondence}

\section{Why the columns have one-letter names}

They are abbreviations of longer names, and the longer names are in a paper from May 2025. The lineage is exact, which is the point of this section: \texttt{o}/\texttt{r} did not appear from nowhere, and reading the ancestor makes the letters legible.

The ancestor is \texttt{CDA.tex}, the Categorical Database Architecture draft. Its schema is three tables with explicit column names. Its \texttt{arrows} table:

\begin{lstlisting}
CREATE TABLE arrows (
  arrow_id SERIAL PRIMARY KEY,
  source_object INTEGER REFERENCES objects(object_id),
  target_object INTEGER REFERENCES objects(object_id),
  arrow_type VARCHAR(100) NOT NULL,
  arrow_data JSONB,
  created_by INTEGER REFERENCES objects(object_id),
  moment TIMESTAMP DEFAULT NOW()
);
\end{lstlisting}

and its \texttt{objects} table:

\begin{lstlisting}
CREATE TABLE objects (
  object_id SERIAL PRIMARY KEY,
  properties JSONB NOT NULL
);
\end{lstlisting}

Line them up:

\begin{center}
\begin{tabular}{llp{5.4cm}}
\toprule
\texttt{CDA.tex} & live & note \\
\midrule
\texttt{objects.object\_id} & \texttt{o.i} & a counter then, a uuidv7 since 2026-07-17 \\
\texttt{objects.properties} & \texttt{o.j} & JSONB, then and now \\
\midrule
\texttt{arrows.arrow\_id} & \texttt{r.i} & lost its primary key on the way \\
\texttt{arrows.source\_object} & \texttt{r.o1} & lost its foreign key \\
\texttt{arrows.target\_object} & \texttt{r.o2} & lost its foreign key \\
\texttt{arrows.arrow\_data} & \texttt{r.d} & \\
\texttt{arrows.arrow\_type} & \texttt{r.d->>'type'} & column folded into the payload \\
\texttt{arrows.created\_by} & \texttt{r.b} & see below \\
\texttt{arrows.moment} & \texttt{r.t} & now a partitioning key \\
\midrule
(none) & \texttt{r.c} & confidence: new \\
\bottomrule
\end{tabular}
\end{center}

\begin{Correspondence}{\texttt{o}/\texttt{r} is CDA's \texttt{objects}/\texttt{arrows} with the names crushed down and \texttt{c} added}{Read the two column lists side by side, the CDA schema section against the tables in \texttt{stem-cell.sql}.}
Seven of the eight \texttt{arrows} columns map onto \texttt{r}. The eighth, \texttt{arrow\_type}, moved inside \texttt{arrow\_data}, where the live code reads it as \texttt{d->>'type'}.
\begin{Check}
grep -n "CREATE TABLE" CDA.tex        -- objects, arrows, evolution
grep -n -A 4 "CREATE TABLE o" stem-cell.sql   -- the live o and r, in full
\end{Check}
The mapping is column-for-column, not thematic. That is worth knowing because it locates precisely what changed and what did not. What did not change: two JSONB payloads carrying everything, identity minted by the database, a timestamp defaulted to \texttt{now()}. Two things changed. One is \texttt{c}, which is the next subsection and is the reason for the book. The other is that the identity stopped being an integer, which took fourteen months and cost a universe, and is Chapter 11.
\end{Correspondence}

\subsection{The addition: \texttt{c}}

\texttt{CDA.tex} has no confidence column, no weight, no uncertainty anywhere at the relation level. Every arrow in it is simply true. Uncertainty appears in that paper exactly once, in Future Work, as ``Uncertainty Propagation: Using observer reliability metrics to propagate measurement uncertainties through complex analysis chains.'' So \textbf{CDA is a certainty database}, and \texttt{c} is the whole difference between it and this one. Everything downstream in this book, the consensus law, calibration, the linear-versus-Born question, the threshold in \texttt{component()}, exists because of one \texttt{REAL NOT NULL DEFAULT 1}.

The correspondence table above hides one substitution worth making explicit. \texttt{arrows.created\_by} pointed at an object: the observer, the person or instrument that made the assertion. The live \texttt{b} is a branch UUID, which is a different thing entirely. Provenance did not move to \texttt{b}; it moved into \texttt{d}. The slot survived and its meaning did not, which is the sort of thing a column-for-column table makes look tidier than it was.

There is also a subtraction inside \texttt{r} itself, recorded in the header of \texttt{schema.sql} and quoted here because it is the clearest statement in the corpus of what \texttt{c} is for:

\begin{quote}
Column w (integer weight) is dropped: multiplicity of evidence is more witness rows; strength of a witness is c. One knob, calibratable.
\end{quote}

Two separate quantities, how many times a thing was seen and how much each sighting is worth, were collapsed into one, on the grounds that the first is already representable as a row count. Whether one knob is enough is a question for Chapter 17, which has to fit it.

\subsection{The subtraction: the third table, twice}

\texttt{CDA.tex} had a third table:

\begin{lstlisting}
CREATE TABLE evolution (
  object_id INTEGER REFERENCES objects(object_id),
  moment TIMESTAMP DEFAULT NOW(),
  state JSONB NOT NULL,
  observer INTEGER REFERENCES objects(object_id),
  PRIMARY KEY (object_id, moment)
);
\end{lstlisting}

An observation: who saw what, of which object, when. \texttt{DESIGN.md}, written fourteen months later and with no visible awareness of the parallel, also specified three tables, and its third was a pair:

\begin{quote}
\texttt{witness(id, ts, o1, o2, evidence\_type, raw\_signal, confidence, provenance, level)} : first-class, retractable \\
\texttt{outcome(id, ts, witness\_id, verdict\{confirmed|retracted\}, cause, provenance)} : the level-2 ledger
\end{quote}

Neither third table survives. Both were absorbed by the same move.

\texttt{evolution} went first: an observation is a relation. There is nothing an \texttt{evolution} row says that a row of \texttt{r} does not say, once you notice that ``observer $O$ saw object $X$ in state $S$ at time $t$'' is an edge from $O$ to $X$ carrying $S$ in its payload and $t$ in its timestamp. \texttt{arrows} was already the general case of \texttt{evolution} and the paper never noticed.

\texttt{outcome} went second, and needed something new: an outcome is a relation pointing at a relation's \texttt{i}. That is only expressible because \texttt{r} has an identity column, which is amendment~1. Retract a witness by writing another row of \texttt{r} whose \texttt{d} carries \texttt{\{"type":"outcome","of":<i>,"verdict":"retracted"\}}. The level-2 ledger is not a separate ledger; it is the same table talking about itself.

\begin{Correspondence}{Reification ate the third table, twice}{Read CDA's \texttt{evolution} table and DESIGN.md's third pair against the live \texttt{\textbackslash d}: there is no third table to find.}
The argument used to be carried by a count. On the universe that ended on 17 July there were 57,994 outcome rows living in \texttt{r}, rows that both prior designs would have given a table of their own, and the number did the arguing. The live instance has zero, because it has zero rows of any kind, so the count now argues nothing and is left here as an honest zero rather than quietly dropped. What survives without the count is the structural half, which was always the real claim: \texttt{\textbackslash dt} returns two tables, and there is no third for an outcome to live in whether or not any outcome exists yet. The absence of the table is checkable today. The population of \texttt{r} is a matter for whenever \labPI{} puts something in it.
\begin{Check}
sed -n '71,72p' DESIGN.md   -- witness and outcome, the third pair
\dt                         -- o and r. No outcome table, no evolution table.
select count(*) from r where d->>'type'='outcome';
--   0 live, out of 0 rows. It was 57994 out of 256619 on the archived universe,
--   backups/ydb/ydb-pre-bigbang-2026-07-17.dump, which is where the count now lives.
\end{Check}
The two eliminations are not the same trick. \texttt{evolution} fell out because it was a special case of a general table that already existed. \texttt{outcome} fell out because arrows became things that can be pointed at. The second is the substantive one, and Chapter 9 argues it is also what makes ``level is a column, not a tier of the software'' true rather than aspirational.
\end{Correspondence}

\section{The three amendments}

They were agreed on 2026-07-03 and were recorded in the header comment of \texttt{schema.sql}; the same three now head \texttt{stem-cell.sql}, which is the file that boots a live instance. They are quoted here in full because the rest of Part II is an audit of whether they hold:

\begin{quote}
1. Arrows are reifiable: r has an identity column i, so witnesses can be referenced by outcomes (the level-2 ledger) via \texttt{d->'of'}.

2. Nothing is ever deleted: \texttt{realize()} records a branch-selection EVENT instead of deleting losing branches. Compression makes keeping everything cheap; there is deliberately NO retention policy.

3. The consensus law is data, not schema: the view below is the rendering of the rule object inserted at bootstrap (linear c-weighting, a seed prior to be replaced by calibration). Regenerate the view when the rule row changes.
\end{quote}

Amendment~1 is the \texttt{i} column of \texttt{r} and it is the whole of the previous section: it is what let the third table go away, and it is what every retraction in Chapter 4 is built on. Amendment~2 is \texttt{b} and \texttt{realize()}; Chapter 8 reports how many times \texttt{realize()} has been called, and the answer is none, on this instance and on the archived one. Amendment~3 is why the \texttt{consensus} view is generated from a row in \texttt{o} rather than being the law itself; Chapter 10 prints that row.

Note the parenthesis in amendment 2: ``Compression makes keeping everything cheap.'' That is the economic argument for the whole design, and it is written in the present tense, in a comment, above a compression policy.

\begin{Open}{Does compression make keeping everything cheap?}{Measure the compressed-to-uncompressed ratio once real data is old enough for the 30-day policy to fire.}
This was unmeasured on 2026-07-16, when the corpus was twelve days old and no chunk was eligible. It is now unmeasured in a stronger sense: the corpus is gone, the live \texttt{r} is empty, and \texttt{timescaledb\_information.chunks} reports no chunks at all for \texttt{r}, so the policy has nothing to run on and did not merely fail to run. The compression job exists and is scheduled, which is the whole of what can be said for it. Every claim in this book about retention being affordable is a prediction, it is a plausible prediction, it is not a measurement, and the comment that asserts it is load-bearing for amendment~2. The sharper form of the question is in Chapter 12: \texttt{compress\_segmentby='b'} segments by a column that held a distinct random UUID on every one of the archived universe's 256,619 rows, which is the worst possible key for a segment-by, and nobody has ever measured what that costs.
\end{Open}

\section{One asymmetry, deferred}

\texttt{o.i} is a \texttt{PRIMARY KEY}. \texttt{r.i} is not. It carries \texttt{NOT NULL} and a non-unique btree index and nothing else.

This is not an oversight. TimescaleDB requires that any unique constraint on a hypertable include the partitioning column, and the partitioning column here is \texttt{t}. A primary key on \texttt{r} would have to be \texttt{(i, t)}, which is not the constraint anyone wants: it permits two rows with the same \texttt{i} at different timestamps, which is precisely the case that matters.

\begin{Correspondence}{\texttt{r.i} is not unique and is not enforced}{Read what index \texttt{r.i} actually carries, then try to add the primary key and read the refusal.}
The refusal is the interesting half. Postgres does not object to the constraint; TimescaleDB does, and it says why.
\begin{Check}
\d r
-- r_i_idx is btree (i), not UNIQUE btree (i), and no primary key is listed
BEGIN;
ALTER TABLE r ADD PRIMARY KEY (i);
-- ERROR: cannot create a unique index without the column "t"
--        (used in partitioning)
ROLLBACK;
\end{Check}
So the identity that every outcome row points at, through the \texttt{(d->>'of')::UUID} cast in \texttt{r\_active}, is a non-unique column with no foreign key referencing it. Nothing in the schema stops two \texttt{r} rows from sharing an \texttt{i}. Whether that is a hole or a footnote is Chapter 11's problem, and Chapter 11 argues it is a hole.
\end{Correspondence}

That is the schema. Two tables, seven columns in the second one, one of which is the reason for the book. What the database does with them is next.
